\documentclass[12pt,technote]{IEEEtran}

\usepackage[a4paper,margin=1in]{geometry}
\usepackage{hyperref}
\usepackage{enumitem}
\usepackage{longtable}
\usepackage{titlesec}

\title{Software Requirements Specification\\
Drop-In FluentValidation-Compatible Validator Library}
\author{Version 1.0}
\date{\today}

\begin{document}
\maketitle

\begin{abstract}
This Software Requirements Specification (SRS) defines the functional and non-functional requirements for a .NET validation library intended to serve as a drop-in replacement for FluentValidation. The system targets .NET 10 and C\# 14, emphasizes minimal reflection, thread safety, Serilog-based diagnostics, and source-generator friendliness. This document follows IEEE SRS conventions and includes a traceability matrix and architectural overview.
\end{abstract}

\tableofcontents
\newpage

\section{Introduction}

\subsection{Purpose}
This SRS defines the requirements for a validation library that is API-compatible with FluentValidation. The primary goals are compatibility, performance, testability, and educational value for understanding validator architecture.

\subsection{Scope}
The system is a .NET 10 class library providing:
\begin{itemize}[noitemsep]
    \item Fluent rule definition
    \item Validation execution (sync and async)
    \item ASP.NET and DI integration
    \item Serilog logging
    \item Minimal reflection
    \item Source-generator friendliness
\end{itemize}

\subsection{Definitions}
\begin{itemize}[noitemsep]
    \item \textbf{Validator} — Class defining validation rules for a model.
    \item \textbf{Rule} — A validation condition applied to a property.
    \item \textbf{ValidationContext} — Encapsulates the instance and metadata for validation.
    \item \textbf{ValidationResult} — Contains validation outcome and failures.
    \item \textbf{Drop-in replacement} — Compatible API enabling substitution for FluentValidation.
\end{itemize}

\subsection{References}
FluentValidation documentation (behavioral reference only).

\subsection{Overview}
This document describes system features, constraints, functional and non-functional requirements, architecture, and traceability.

\section{Overall Description}

\subsection{Product Perspective}
The validator is a standalone .NET library intended to mimic FluentValidation’s API and behavior. It integrates with ASP.NET model validation and DI.

\subsection{Product Functions}
\begin{itemize}[noitemsep]
    \item Fluent rule definition (\texttt{RuleFor}, \texttt{NotEmpty}, etc.)
    \item Validation execution (\texttt{Validate}, \texttt{ValidateAsync})
    \item Error reporting (\texttt{ValidationResult}, \texttt{ValidationFailure})
    \item ASP.NET integration
    \item Logging via Serilog
    \item Extensibility via custom validators
\end{itemize}

\subsection{User Characteristics}
Users are .NET developers familiar with FluentValidation, DI, and ASP.NET.

\subsection{Constraints}
\begin{itemize}[noitemsep]
    \item Must run on .NET 10 and C\# 14
    \item Must minimize reflection
    \item Must be thread-safe
    \item Must integrate with Serilog
    \item Must be source-generator friendly
\end{itemize}

\subsection{Assumptions and Dependencies}
\begin{itemize}[noitemsep]
    \item FluentValidation API is stable enough to model
    \item Serilog is configured by the host application
    \item A DI container is available
\end{itemize}

\section{Functional Requirements}

\subsection{Rule Definition and Composition}
\begin{itemize}[noitemsep]
    \item \textbf{FR-1} The system shall provide \texttt{AbstractValidator<T>} with \texttt{RuleFor} methods.
    \item \textbf{FR-2} The system shall support built-in validators: \texttt{NotNull}, \texttt{NotEmpty}, \texttt{Length}, \texttt{Matches}, \texttt{EmailAddress}, etc.
    \item \textbf{FR-3} The system shall support chaining of validators.
    \item \textbf{FR-4} The system shall support conditional rules (\texttt{When}, \texttt{Unless}).
    \item \textbf{FR-5} The system shall support rule sets.
\end{itemize}

\subsection{Custom Validators}
\begin{itemize}[noitemsep]
    \item \textbf{FR-6} The system shall support custom validator classes.
    \item \textbf{FR-7} The system shall support predicate-based validators (\texttt{Must}).
    \item \textbf{FR-8} Custom validators shall support custom error messages and error codes.
\end{itemize}

\subsection{Validation Execution}
\begin{itemize}[noitemsep]
    \item \textbf{FR-9} The system shall provide \texttt{Validate} and \texttt{ValidateAsync}.
    \item \textbf{FR-10} The system shall evaluate all rules by default.
    \item \textbf{FR-11} The system shall support cascade modes.
\end{itemize}

\subsection{Error Messages and Localization}
\begin{itemize}[noitemsep]
    \item \textbf{FR-12} The system shall provide default error messages.
    \item \textbf{FR-13} The system shall allow overriding messages via \texttt{WithMessage()}.
    \item \textbf{FR-14} The system shall support localization.
\end{itemize}

\subsection{Configuration}
\begin{itemize}[noitemsep]
    \item \textbf{FR-15} The system shall provide global configuration for cascade mode and message providers.
    \item \textbf{FR-16} Configuration shall be thread-safe.
\end{itemize}

\section{Non-Functional Requirements}

\subsection{Performance and Reflection}
\begin{itemize}[noitemsep]
    \item \textbf{NFR-1} Reflection usage shall be minimized.
    \item \textbf{NFR-2} Reflection results shall be cached.
    \item \textbf{NFR-3} The system shall use compiled delegates where possible.
\end{itemize}

\subsection{Thread Safety}
\begin{itemize}[noitemsep]
    \item \textbf{NFR-4} Validators shall be safe for concurrent use.
    \item \textbf{NFR-5} Shared state shall be immutable or synchronized.
\end{itemize}

\subsection{Testability}
\begin{itemize}[noitemsep]
    \item \textbf{NFR-6} The system shall exhibit deterministic behavior.
    \item \textbf{NFR-7} The system shall avoid hidden global state.
\end{itemize}

\subsection{Source-Generator Friendliness}
\begin{itemize}[noitemsep]
    \item \textbf{NFR-8} Rule definition shall be separable from execution.
    \item \textbf{NFR-9} APIs shall exist for precompiled rule metadata.
\end{itemize}

\subsection{Reliability}
\begin{itemize}[noitemsep]
    \item \textbf{NFR-10} Validation failures shall not throw exceptions.
    \item \textbf{NFR-11} Internal errors shall be logged via Serilog.
\end{itemize}

\section{Architecture Overview}

\subsection{Core Components}
\begin{itemize}[noitemsep]
    \item \textbf{Validator Engine} — Executes rules and produces results.
    \item \textbf{Rule Builder} — Fluent API for constructing rules.
    \item \textbf{Property Validators} — Built-in and custom validators.
    \item \textbf{Execution Pipeline} — Applies rules with cascade logic.
    \item \textbf{Integration Layer} — ASP.NET and DI support.
    \item \textbf{Logging Layer} — Serilog instrumentation.
\end{itemize}

\subsection{Data Flow}
\begin{enumerate}[noitemsep]
    \item User defines validator via fluent API.
    \item Validator compiles rules into delegates.
    \item Validation engine executes rules.
    \item Failures are collected into \texttt{ValidationResult}.
    \item Serilog logs diagnostics.
\end{enumerate}

\section{Recommended Project Structure}

\begin{verbatim}
src/
  Core/
    AbstractValidator.cs
    IValidator.cs
    ValidationResult.cs
    ValidationFailure.cs
  Rules/
    RuleBuilder.cs
    PropertyValidators/
  Execution/
    ValidationContext.cs
    ValidatorEngine.cs
  Integration/
    AspNet/
    DependencyInjection/
  Logging/
    SerilogAdapter.cs

tests/
  CoreTests/
  RuleTests/
  IntegrationTests/
\end{verbatim}

\section{Traceability Matrix}

\begin{longtable}{|p{3cm}|p{10cm}|}
\hline
\textbf{Requirement} & \textbf{Implementation Component} \\
\hline
FR-1--FR-5 & RuleBuilder, AbstractValidator \\
FR-6--FR-8 & CustomValidator API \\
FR-9--FR-11 & ValidatorEngine \\
FR-12--FR-14 & MessageProvider, Localization \\
FR-15--FR-16 & GlobalConfiguration \\
NFR-1--NFR-3 & ReflectionCache, DelegateCompiler \\
NFR-4--NFR-5 & Thread-safe configuration and caching \\
NFR-6--NFR-7 & Test harness, deterministic pipeline \\
NFR-8--NFR-9 & Source generator APIs \\
NFR-10--NFR-11 & SerilogAdapter, exception handling \\
\hline
\end{longtable}

\section{Glossary}
\begin{itemize}[noitemsep]
    \item \textbf{AOT} — Ahead-of-time compilation.
    \item \textbf{DI} — Dependency Injection.
    \item \textbf{Pipeline} — Sequence of rule evaluations.
\end{itemize}

\end{document}
